\documentclass[10pt,a4paper]{amsart}
\usepackage[margin=19mm]{geometry}
\usepackage{amsmath,amssymb,mathtools}
\usepackage[scr=rsfs,cal=euler]{mathalfa}
\usepackage{xfrac}
\usepackage[unicode,hidelinks,bookmarks=false]{hyperref}
\newcommand{\ie}{i.e.\ }
\newcommand{\cf}{cf.\ }
\newcommand{\resp}{respectively\ }
\newcommand{\Subsection}{\subsection}
\providecommand{\nobreakdash}{\nobreak\hbox{-}}
\setlength{\emergencystretch}{2em}
\multlinegap0pt
\usepackage{amssymb}
\usepackage{mathtools}
\usepackage[shortlabels]{enumitem}
\newtheorem{lemma}{Lemma}
\title{\textbf{A Symplectic Proof of the Quantum Singleton Bound}}
\author{Frederick Dehmel, Shilun Li}
\date{Source: arXiv:2602.20186v3 (CC BY 4.0)}
\begin{document}
\maketitle

\noindent\emph{Source and licence.} \textbf{A Symplectic Proof of the Quantum Singleton Bound}. Version arXiv:2602.20186v3 (CC BY 4.0), distributed by arXiv under the Creative Commons Attribution 4.0 International licence, http://creativecommons.org/licenses/by/4.0/, as recorded in the arXiv licence metadata for this exact version and shown on the abstract page https://arxiv.org/abs/2602.20186v3. Full licence notice: \texttt{licenses/arxiv-2602.20186-CC-BY-4.0.txt}.\par\smallskip

\noindent\textit{Editorial scope.} Counted unit: the article's lemma lem:two-disjoint (source0.tex lines 434--440). The article's complete original proof of it is delivered with it (source0.tex lines 442--475, one \texttt{proof} environment of 34 lines). No mathematics is written, repaired or re-derived: the statement and the proof are the article's own lines, in the article's own notation, and no retained line is substituted or re-flowed.  Retained uncounted as the source-local context the proof needs: lines 351--356.  Nothing was omitted from any retained span, so the package carries no omission record.  No retained line contains a citation, so the wrapper carries no bibliography and no bibliography entry.  No retained line contains a figure environment, an image-inclusion command or drawing code, so no image is delivered and the package declares no asset dependency.  Editorial formatting only, in addition: an AMS \texttt{amsart} preamble that restates the article's own declarations and macros used by the retained lines, namely the environments \texttt{lemma} on separate counters, together with the standard packages those lines need.  The retained environments print in this document as Lemma 1 in order of appearance, whereas the article prints the same items with the numbering of its own full document.  The complete native source of the article as distributed in the arXiv e-print for this exact version is bundled unchanged as \texttt{sources/195166/source0.tex}.
\begin{lemma}[Cleaning dimension identity]\label{lem:cleaning}
For any $M \subseteq [n]$,
\[
  g(M) + g(M^c) = 2k.
\]
\end{lemma}

\begin{lemma}[Two disjoint correctable sets bound logical dimension]\label{lem:two-disjoint}
Partition the physical qudits into $A\, B\, C$ with $A \cap B = \varnothing$ and
$C = [n] \setminus (A \cup B)$.  If erasures of both $A$ and $B$ are correctable, then
\[
  k \le |C|.
\]
\end{lemma}

\begin{proof}
Set $D := A^c = B \cup C$.  Since $A$ is correctable we have $g(A) = 0$, hence by the
cleaning dimension identity (Lemma~\ref{lem:cleaning}),
\[
  g(D) = g(A^c) = 2k.
\]
By definition,
\[
  g(D) = \dim\Bigl( (S^\perp \cap V_D) / (S \cap V_D) \Bigr),
\]
so the quotient space $L := (S^\perp \cap V_D) / (S \cap V_D)$ has dimension $\dim(L) = 2k$.

Consider the restriction map $r_C : V_D \to V_C$ (zeroing coordinates outside $C$).  Let
$W := S \cap V_D$.  Then $r_C(W) \le V_C$ is a subspace, and $r_C$ induces a linear map
\[
  \bar{r}_C : L \longrightarrow V_C / r_C(W),
  \qquad [v] \longmapsto r_C(v) \bmod r_C(W).
\]
This is well-defined because replacing $v$ by $v + w$ with $w \in W$ changes $r_C(v)$ by an
element of $r_C(W)$.

We claim $\bar{r}_C$ is injective.  Suppose $\bar{r}_C([v]) = 0$.  Then
$r_C(v) \in r_C(W)$, so there exists $w \in W$ with $r_C(v) = r_C(w)$.  Hence
$r_C(v - w) = 0$, so $v - w$ is supported entirely on $B$, i.e.,
$v - w \in V_B$.  Also $v - w \in S^\perp$ (since $v \in S^\perp$ and $w \in S$).  Thus
$v - w \in S^\perp \cap V_B$, and because $B$ is correctable we get $v - w \in S$.  Since
$w \in S$ as well, it follows that $v \in S$, hence $[v] = 0$ in $L$.

Therefore,
\[
  2k = \dim(L) \le \dim\bigl( V_C / r_C(W) \bigr) \le \dim(V_C) = 2|C|,
\]
so $k \le |C|$.
\end{proof}

\end{document}
